% General Article Template — Texforge Capabilities Demo
\documentclass[a4paper,12pt]{article}
\usepackage[utf8]{inputenc}
\usepackage[T1]{fontenc}
\usepackage[spanish]{babel}
\usepackage{amsmath,amssymb}
\usepackage{float}
\usepackage{graphicx}
\usepackage{hyperref}
\usepackage{geometry}
\usepackage{fancyhdr}
\usepackage{booktabs}
\usepackage{enumitem}
\usepackage{listings}
\usepackage{xcolor}

\geometry{margin=2.5cm}

\pagestyle{fancy}
\fancyhf{}
\cfoot{\thepage}

\lstset{
  basicstyle=\ttfamily\small,
  breaklines=true,
  frame=single,
  backgroundcolor=\color{gray!10}
}

\title{Demostración de Capacidades de Texforge}
\author{Jheison M}
\date{\today}

\begin{document}

  \maketitle

  \begin{abstract}
    Este documento demuestra las capacidades principales de \textbf{texforge},
    un compilador LaTeX que utiliza tectonic como motor interno, eliminando la
    necesidad de instalar TeX Live o MiKTeX. Se presentan diagramas embebidos
    (Mermaid, Graphviz y D2), citas bibliográficas, referencias cruzadas,
    matemáticas, listados de código, tablas y más.
  \end{abstract}

  \tableofcontents
  \newpage

  % -------------------------------------------------------
  \section{Introducción}\label{sec:intro}

  Texforge simplifica el flujo de trabajo académico al compilar documentos
  LaTeX directamente a PDF sin distribuciones externas~\cite{example}. En
  esta sección se describen los antecedentes y motivaciones del proyecto.

  \subsection{Antecedentes}

  La edición de documentos académicos en \LaTeX{} requiere tradicionalmente
  instalaciones pesadas como TeX Live o MiKTeX. Texforge elimina esa
  barrera~\cite{example2}, ofreciendo una experiencia de línea de comandos
  ligera y moderna.

  \subsection{Objetivos}

  Los objetivos de este documento son:

  \begin{enumerate}[label=(\roman*)]
    \item Demostrar diagramas embebidos en tres formatos distintos.
    \item Mostrar la integración bibliográfica con BibTeX.
    \item Ejemplificar referencias cruzadas entre secciones.
    \item Presentar ecuaciones matemáticas complejas.
    \item Incluir listados de código fuente.
  \end{enumerate}

  % -------------------------------------------------------
  \section{Diagramas Embebidos}\label{sec:diagramas}

  Texforge renderiza diagramas directamente desde el código fuente LaTeX
  sin herramientas externas como Inkscape o Node.js.

  \subsection{Mermaid}

  El siguiente diagrama de secuencia muestra la interacción entre los
  componentes de Texforge:

  \begin{mermaid}[caption=Secuencia de compilación de Texforge, height=10cm, pos=t]
    sequenceDiagram
    actor User as Usuario
    participant CLI as texforge
    participant PRE as Pre-procesador
    participant TEC as tectonic

    User->>CLI: texforge build
    CLI->>PRE: parsea .tex y detecta diagramas
    loop Diagramas embebidos
    PRE->>CLI: Mermaid / Graphviz / D2
    CLI-->>PRE: imagen PNG generada
    end
    CLI->>TEC: main.tex + imágenes
    TEC-->>User: PDF final
  \end{mermaid}

  \subsection{Graphviz / DOT}

  También se pueden crear diagramas directed-Graph con Graphviz:

  \begin{graphviz}[caption=Arquitectura de Texforge, width=0.8\linewidth, pos=t]
    digraph G {
      rankdir=LR
      node [shape=box style="rounded,filled" fillcolor="#E8F4FD"]

      CLI [label="CLI\ntexforge"]
      Parser [label="Parser\nLaTeX"]
      MermaidEngine [label="Mermaid\nRenderer"]
      GraphvizEngine [label="Graphviz\nRenderer"]
      D2Engine [label="D2\nRenderer"]
      Tectonic [label="tectonic\nEngine"]
      PDF [label="Output\nPDF"]

      CLI -> Parser
      Parser -> MermaidEngine
      Parser -> GraphvizEngine
      Parser -> D2Engine
      MermaidEngine -> Tectonic
      GraphvizEngine -> Tectonic
      D2Engine -> Tectonic
      Tectonic -> PDF
    }
  \end{graphviz}

  \subsection{D2}

  El diagrama D2 es otro formato soportado nativamente:

  \begin{d2}[caption=Pipeline de renderizado, height=14cm, pos=t]
    source: "Código LaTeX" {
      shape: document
    }
    preprocessor: "Pre-procesador" {
      shape: hexagon
    }
    engines: "Engines de diagramas" {
      shape: package
    }
    pdf: "PDF final" {
      shape: page
    }

    source -> preprocessor: "lex & parse"
    preprocessor -> engines: "diagramas detectados"
    engines -> pdf: "tectonic"
  \end{d2}

  \subsection{Estilos de Diagrama (\texttt{style})}\label{sec:estilos}

  Los tres entornos aceptan también un atributo \texttt{style}, combinable
  con \texttt{width}, \texttt{pos} y \texttt{caption}. En vez de describir
  cada preset con adjetivos, el mismo diagrama mínimo se repite a
  continuación una vez por preset, para que la diferencia se vea
  directamente:

  \begin{mermaid}[style=editorial, width=0.55\linewidth, caption=Preset \texttt{editorial}: paleta restringida, un solo acento, sin sombras, pos=H]
    flowchart LR
      A[Entrada] --> B[Proceso] --> C[Salida]
  \end{mermaid}

  \begin{mermaid}[style=mono, width=0.55\linewidth, caption=Preset \texttt{mono}: escala de grises, pos=H]
    flowchart LR
      A[Entrada] --> B[Proceso] --> C[Salida]
  \end{mermaid}

  \begin{mermaid}[style=technical, width=0.55\linewidth, caption=Preset \texttt{technical}: trazo uniforme, sin relleno, etiquetas monoespaciadas, pos=H]
    flowchart LR
      A[Entrada] --> B[Proceso] --> C[Salida]
  \end{mermaid}

  El preset \texttt{editorial} sirve para documentos que se leen en pantalla
  o se imprimen a color y buscan un acabado sobrio. El preset
  \texttt{technical} conviene a diagramas que funcionan como planos o
  esquemas de referencia, con peso de trazo uniforme. El preset
  \texttt{mono} tiene una razón menos obvia: existe para documentos que se
  imprimen en blanco y negro, donde un diagrama a color puede perder
  significado, porque dos colores que se distinguen en pantalla pueden
  reducirse al mismo gris al imprimirse sin color; \texttt{mono} evita ese
  problema por construcción, ya que nunca depende del color para
  distinguir nada. Omitir \texttt{style} equivale a \texttt{style=default},
  el aspecto propio de cada motor, como en los tres diagramas anteriores.

  También puede fijarse un preset por defecto para todo el documento en
  \texttt{project.toml}, bajo \texttt{[diagrams]} con \texttt{style = "..."};
  el atributo \texttt{style} de un entorno concreto, cuando está presente,
  sobrescribe ese valor por defecto.

  % -------------------------------------------------------
  \section{Matemáticas}\label{sec:matematicas}

  \subsection{Ecuaciones en línea}

  La fórmula cuadrática se expresa como $x = \frac{-b \pm \sqrt{b^2 - 4ac}}{2a}$.

  \subsection{Ecuaciones numeradas}

  \begin{equation}\label{eq:gauss}
    \int_{-\infty}^{\infty} e^{-x^2}\,dx = \sqrt{\pi}
  \end{equation}

  \begin{equation}\label{eq:euler}
    e^{i\pi} + 1 = 0
  \end{equation}

  Como se demostró en la Ecuación~\ref{eq:gauss}, la integral de Gauss
  tiene un valor exacto conocido.

  \subsection{Matrices}

  \begin{equation}\label{eq:matriz}
    \mathbf{A} = \begin{pmatrix}
      a_{11} & a_{12} & \cdots & a_{1n} \\
      a_{21} & a_{22} & \cdots & a_{2n} \\
      \vdots & \vdots & \ddots & \vdots \\
      a_{n1} & a_{n2} & \cdots & a_{nn}
    \end{pmatrix}
  \end{equation}

  % -------------------------------------------------------
  \section{Citas Bibliográficas}\label{sec:bib}

  Texforge gestiona automáticamente las referencias BibTeX. A continuación
  se presentan tres tipos de cita:

  \begin{itemize}
    \item Artículos de revista~\cite{example}
    \item Libros~\cite{example2}
    \item Recursos en línea~\cite{example3}
  \end{itemize}

  La sección~\ref{sec:matematicas} muestra el uso de referencias cruzadas,
  mientras que la sección~\ref{sec:diagramas} describe los formatos de
  diagramas disponibles.

  % -------------------------------------------------------
  \section{Listados de Código}\label{sec:codigo}

  \subsection{Python}

  \begin{lstlisting}[language=Python, caption={Ejemplo de función en Python}]
import numpy as np

def integrar_gauss(n_puntos=1000):
    """Aproxima la integral de Gauss por cuadratura."""
    x = np.linspace(-10, 10, n_puntos)
    dx = x[1] - x[0]
    y = np.exp(-x**2)
    return np.sum(y) * dx

if __name__ == "__main__":
    resultado = integrar_gauss()
    print(f"Integral ≈ {resultado:.6f}")
  \end{lstlisting}

  \subsection{LaTeX}

  \begin{lstlisting}[language={[LaTeX]TeX}, caption={Estructura básica de Texforge}]
\documentclass{article}
\usepackage{graphicx}

\begin{document}
\maketitle

\begin{mermaid}
flowchart LR
  A --> B --> C
\end{mermaid}

\bibliographystyle{plain}
\bibliography{refs}
\end{document}
  \end{lstlisting}

  % -------------------------------------------------------
  \section{Tablas}\label{sec:tablas}

  \begin{table}[H]
    \centering
    \caption{Comparación de motores de renderizado de diagramas}
    \label{tab:motores}
    \begin{tabular}{@{}llcc@{}}
      \toprule
      \textbf{Formato} & \textbf{Motor} & \textbf{Externo?} & \textbf{Velocidad} \\
      \midrule
      Mermaid    & Rust puro         & No  & Rápido   \\
      Graphviz   & layout-rs         & No  & Rápido   \\
      D2         & d2-little         & No  & Rápido   \\
      \bottomrule
    \end{tabular}
  \end{table}

  Como se muestra en la Tabla~\ref{tab:motores}, todos los motores son
  embebidos y no requieren binarios externos.

  % -------------------------------------------------------
  \section{Comandos de Texforge}\label{sec:comandos}

  \begin{table}[H]
    \centering
    \caption{Resumen de comandos principales}
    \label{tab:comandos}
    \begin{tabular}{@{}lp{8cm}@{}}
      \toprule
      \textbf{Comando} & \textbf{Descripción} \\
      \midrule
      \texttt{texforge new}    & Crea un proyecto nuevo con un template. \\
      \texttt{texforge init}   & Wizard interactivo para proyectos existentes. \\
      \texttt{texforge build}  & Compila \texttt{main.tex} a PDF en \texttt{build/}. \\
      \texttt{texforge check}  & Linter estático: valida archivos, citas y labels. \\
      \texttt{texforge fmt}    & Formatea archivos \texttt{.tex}. \\
      \texttt{texforge clean}  & Elimina el directorio \texttt{build/}. \\
      \texttt{texforge template} & Gestiona templates instalados o remotos. \\
      \bottomrule
    \end{tabular}
  \end{table}

  % -------------------------------------------------------
  \section{Conclusión}\label{sec:conclusion}

  Texforge demuestra ser una herramienta completa para la edición de
  documentos académicos en \LaTeX{}. Sus ventajas principales son:

  \begin{itemize}
    \item \textbf{Sin dependencias externas:} no requiere TeX Live ni MiKTeX.
    \item \textbf{Diagramas embebidos:} Mermaid, Graphviz y D2 renderizados en Rust puro.
    \item \textbf{Linter integrado:} detección temprana de errores con \texttt{texforge check}.
    \item \textbf{Templates:} soporte para APA, IEEE, cartas y documentos generales.
  \end{itemize}

  Como se demostró en este documento, la combinación de diagramas
  (Sección~\ref{sec:diagramas}), ecuaciones (Sección~\ref{sec:matematicas}),
  citas (Sección~\ref{sec:bib}) y código (Sección~\ref{sec:codigo}) permite
  crear publicaciones académicas de alta calidad con un mínimo esfuerzo.

  \bibliographystyle{plain}
  \bibliography{bib/references}

\end{document}
